% Weekly meeting deck. Build with pdflatex twice.
% Evidence: docs/experiments/biological_first_results_2026-10-15.md;
% docs/experiments/predictor_reward_study_2026-11-03.md;
% docs/verification/predictor-reward-results-2026-11-03.json.
% No slide source footers or section-navigation headers, per presentation preference.
\documentclass[10pt,aspectratio=169]{beamer}
\geometry{paperwidth=180mm,paperheight=101.25mm}
\usepackage[T1]{fontenc}
\usepackage{lmodern}
\usepackage{booktabs}
\usepackage{ragged2e}
\setlength{\RaggedRightRightskip}{0pt plus 1fil}

% Match the earlier project-progress deck: shared text/title margins,
% top-aligned frames, sans-serif type and bottom-aligned takeaways.
\definecolor{navy}{HTML}{233B53}
\definecolor{teal}{HTML}{207B80}
\definecolor{ink}{HTML}{293541}
\definecolor{muted}{HTML}{5D6B76}
\definecolor{pale}{HTML}{EDF4F5}
\setbeamercolor{normal text}{fg=ink,bg=white}
\setbeamercolor{structure}{fg=teal}
\setbeamercolor{frametitle}{fg=navy}
\setbeamercolor{block title}{fg=navy,bg=pale}
\setbeamercolor{block body}{fg=ink,bg=pale}
\setbeamercolor{takeaway}{fg=navy,bg=pale}
\setbeamerfont{frametitle}{size=\Large,series=\bfseries}
\setbeamerfont{block title}{series=\bfseries}
\setbeamertemplate{navigation symbols}{}
\setbeamertemplate{headline}{}
\setbeamertemplate{itemize items}[circle]
\setbeamersize{text margin left=10mm,text margin right=10mm}
\setbeamertemplate{frametitle}{%
  \vspace{2mm}
  \usebeamerfont{frametitle}\insertframetitle\par
  \vspace{1mm}
}
\setbeamertemplate{footline}{%
  \hfill{\fontsize{8}{9}\selectfont\color{muted}
  \insertframenumber\,/\,\inserttotalframenumber}\hspace{10mm}\vspace{3mm}
}
\setlength{\tabcolsep}{5pt}
\renewcommand{\arraystretch}{1.17}
\setlength{\parskip}{3pt}
\hyphenpenalty=10000
\exhyphenpenalty=10000
\newcommand{\takeaway}[1]{%
  \par\vspace{2mm}\vfill
  \begin{beamercolorbox}[wd=\textwidth,sep=2mm]{takeaway}
    \small\RaggedRight\leftskip=2mm\rightskip=2mm plus 1fil #1
  \end{beamercolorbox}
}
\title{Stability guidance: weekly update}
\date{3 November 2026}
\begin{document}

\begin{frame}[t]{Where we started: the guidance system works}
\RaggedRight
\textbf{Goal:} steer a frozen ProGen2 model toward stable proteins by training
an additional GFlowNet guidance head.

\begin{columns}[T,onlytextwidth]
\column{.31\textwidth}
\RaggedRight
\begin{block}{1. ProGen2}
Provides sequence probabilities and prefix representations. Its weights stay frozen.
\end{block}
\column{.31\textwidth}
\RaggedRight
\begin{block}{2. Stability predictor}
Scores the whole protein. Higher predicted stability increases its target sampling weight.
\end{block}
\column{.31\textwidth}
\RaggedRight
\begin{block}{3. Guidance head}
Learns token-level corrections using trajectory balance, to sample the target distribution.
\end{block}
\end{columns}
\vspace{2mm}
\begin{itemize}
\item ProGen2 and a standalone GFlowNet were working independently.
\item Integration passed a 256-sequence synthetic test: distribution distance
about 0.04--0.05 across three seeds (zero means a perfect match).
\end{itemize}
\end{frame}

\begin{frame}[t]{The first stability result: accurate sampling, no gain}
\RaggedRight
\textbf{Task:} Mega-scale stability data for one 72-residue protein, 1O6X.
950 training sequences, 204 validation sequences, 204 evaluation variants.
Evaluation variants are single amino-acid substitutions.

\begin{center}
\renewcommand{\arraystretch}{1.2}
\begin{tabular}{lrr}
\toprule
 & \multicolumn{2}{c}{Expected stability (kcal/mol)}\\
Distribution & Predicted & Measured\\
\midrule
Constrained ProGen & 3.894 & 4.280\\
Exact predicted-reward target & 3.973 & 4.270\\
Trained GFlowNet (three-seed mean) & 3.973 & 4.270\\
\bottomrule
\end{tabular}
\end{center}
Probability-weighted averages over the same 204 variants; higher is better.
``Exact'' means probabilities are computed directly, without sampling.

\takeaway{The GFlowNet learned its target accurately. The target itself did not
improve measured stability, so we investigated the predictor and reward.}
\end{frame}

\begin{frame}[t]{This week: test the predictor and inspect the reward}
\RaggedRight
\begin{columns}[T,onlytextwidth]
\column{.49\textwidth}
\RaggedRight
\begin{block}{A fair predictor comparison}
\begin{itemize}
\item Constant training-mean control.
\item Ridge: a regularized linear model.
\item MLP: a small neural network.
\end{itemize}
\end{block}
Same frozen ProGen features and partitions. Train on \textbf{40, 190 or 950}
examples; three seeds with matched subsets. Each subset includes the reference.

Choose predictor settings using validation error only.
\column{.47\textwidth}
\RaggedRight
\begin{block}{A focused reward audit}
\begin{itemize}
\item Reconstruct the original score and check normalization.
\item Test guidance strength $\beta=0,\,0.5,\,1,\,2$.
\item Compare the original reward with a reward capped at predicted reference stability.
\end{itemize}
\end{block}
$\beta$: guidance strength. Zero recovers ProGen; one is the original setting.
\end{columns}
\takeaway{Evaluate every target directly before another GFlowNet fit.
The original reward arithmetic checked out.}
\end{frame}

\begin{frame}[t]{The linear predictor gives a better target on this split}
\RaggedRight
\begin{columns}[T,onlytextwidth]
\column{.43\textwidth}
\RaggedRight
\textbf{More data improves prediction}
\begin{center}
\renewcommand{\arraystretch}{1.22}
\begin{tabular}{rrr}
\toprule
Train size & Ridge & MLP\\
\midrule
40 & 1.001 & 0.982\\
190 & 0.738 & 0.712\\
950 & 0.524 & 0.529\\
\bottomrule
\end{tabular}
\end{center}
Evaluation root mean squared error, in kcal/mol; lower is better.
Three-seed means. Constant control: about 1.16.
\column{.53\textwidth}
\RaggedRight
\textbf{Full-data targets, original strength}
\begin{center}
\renewcommand{\arraystretch}{1.22}
\begin{tabular}{lrr}
\toprule
 & Measured & Change\\
\midrule
ProGen & 4.280 & ---\\
Ridge target & \textbf{4.320} & \textbf{+0.040}\\
MLP target & 4.261 & $-0.019$\\
\bottomrule
\end{tabular}
\end{center}
Expected stability in kcal/mol, with $\beta=1$ and 950 training examples.
The MLP loses measured utility in all three seeds.
\end{columns}
\takeaway{Similar overall prediction errors can produce different guidance outcomes.
These are exact target results; no new GFlowNet was trained this week.}
\end{frame}

\begin{frame}[t]{Why: errors on likely variants matter most}
\RaggedRight
ProGen puts \textbf{91.4\% of its probability on five variants}.
Errors in this small group can dominate the effect of guidance.

\begin{center}
\renewcommand{\arraystretch}{1.25}
\begin{tabular}{lrrrr}
\toprule
Variant & ProGen probability & Measured & Ridge prediction & MLP prediction\\
\midrule
V1R & 33.9\% & 4.438 & 4.684 & 4.183\\
G2I & 10.4\% & 3.830 & 3.904 & 4.816\\
D3G & 4.4\% & 4.419 & 4.542 & 3.157\\
\bottomrule
\end{tabular}
\end{center}
Full-data seed-0 models; kcal/mol. G2I: glycine at position 2 replaced by
isoleucine (same naming convention for other variants).

\begin{itemize}
\item The neural model ranks G2I above both more-stable alternatives.
\item All three neural models beat ridge on validation error, but produce worse targets.
\end{itemize}
\takeaway{Evaluate prediction errors where ProGen puts its probability.}
\end{frame}

\begin{frame}[t]{What we learned, and the next focused step}
\RaggedRight
\begin{itemize}
\setlength{\itemsep}{4pt}
\item \textbf{The linear baseline is competitive.} At full data, its target gains
about 0.040 kcal/mol over ProGen on this split.
\item \textbf{Reward capping is not a complete fix.} For the neural predictor,
it slightly raises mean stability, but lowers the chance of reaching reference-level stability.
\item \textbf{The limitation is still local evidence.} One protein and an already
inspected evaluation set; repeated full-data ridge fits are not independent replications.
\end{itemize}
\vspace{2mm}
\begin{block}{Proposed next step}
Fix ridge and the original guidance strength; check additional predefined
splits of this protein's data against the neural comparator. Measure overall
and ProGen-weighted prediction errors, plus exact target quality.

If the gain persists, test whether GFlowNet guidance realizes it. Keep trajectory
balance for now: target matching has already worked.
\end{block}
\end{frame}
\end{document}
